% H200 synthesis: block-preconditioned and multigrid solvers on one H200 against PyPardiso LU on
% a genoa CPU node. Figures and tables: scripts/reports/make_adr_h200_comparison.py
% (grids and envelopes drawn with the Blackwell synthesis' own functions); no solver execution.
\section{H200 synthesis: iterative solvers on the GPU versus PyPardiso LU on a CPU node}
\label{sec:adr-h200-synthesis}

This section repeats the synthesis of the Blackwell study for the H200: every solver family on
one NVIDIA H200 NVL with the tuned kernels (\texttt{campaigns/adr\_h200\_tuned}, 28--29 October
2026), against multithreaded PyPardiso LU on a CPU node of the same cluster (genoa partition, two
AMD EPYC 9655, 192 cores, 754~GiB; PARDISO v2 campaign, jobs 24478 and 24482,
\texttt{campaigns/adr\_pardiso/records/unified\_l5\_pardiso\_v2}). The LU did not run on the H200
node's own host CPUs (two 64-core sockets). LU and the iterative solvers see the same discrete
systems (mesh, degree, case and coefficients, joined by the plan key), assembled twice: the LU
campaign factorized the stock-kernel exports of the unified device campaign, and the tuned
campaign assembled its own. The two assemblies hash differently: on four checked systems the
mesh topology is identical and the blocks and right-hand side agree to $5\cdot10^{-14}$ to
$6\cdot10^{-12}$ relative.

\paragraph{Problems.} Every case is the stationary conservative problem
\begin{equation}
 -\nabla\cdot(K\nabla u)+\nabla\cdot(\boldsymbol\beta u)+c\,u=f\ \text{in }\Omega,\qquad
 u=u_\star\ \text{on }\partial\Omega,
 \label{eq:adr-h200-problem}
\end{equation}
with $f$ manufactured from the exact solution $u_\star$, discretized by HDG of degree $p$ and
condensed to the trace. Table~\ref{tab:adr-h200-cases} lists the coefficients, using
\begin{itemize}
\item \textbf{domains} $\Omega_\square=[-1,1]^2$, the five-lobed annulus
  $\Omega_5=\{0.58<r<1+0.35\cos5\theta\}$ and the nine-lobed annulus
  $\Omega_9=\{0.63<r<R_9\}$, $R_9=1+0.35\cos9\theta$;
\item \textbf{solutions} $u_s=\sin\pi x\sin\pi y$, $u_q=1+x^2+2y^2$,
  $u_o=\sin6\pi x\sin5\pi y+0.35\sin11\pi x\sin9\pi y$ and
  $u_9=\sin2\pi\chi+\tfrac14u_o$, with $\rho=(r-0.63)/(R_9-0.63)$ and
  $\chi=\rho+0.05\sin5\theta\sin2\pi\rho$;
\item \textbf{velocities} $\boldsymbol\beta_0=(1,\tfrac12)$, $\boldsymbol\beta_v=(1+x,-\tfrac12+y)$
  ($\nabla\cdot\boldsymbol\beta_v=2$), the cellular flows
  $\boldsymbol\beta_\omega=(1+4\sin\omega\pi x\cos\omega\pi y,\ \tfrac12-4\cos\omega\pi x\sin\omega\pi y)$,
  and the through-flow $\boldsymbol\beta_\perp=\tfrac{50}{1.8}[1+0.8\sin(13\pi\,e_d\cdot\boldsymbol x)]\,e_t$
  with $e_d=(\cos\tfrac\pi7,\sin\tfrac\pi7)$, $e_t=e_d^\perp$;
\item \textbf{tensors} $K_{a,b}=R(\tfrac\pi6)\operatorname{diag}(a,b)R(\tfrac\pi6)^T$ and
  $K_\perp=e_de_d^T+10^{-6}e_te_t^T$.
\end{itemize}

\begin{table}[ht]
\centering\footnotesize
\setlength{\tabcolsep}{4pt}
\begin{tabular}{llllllll}
\toprule
class & key & $\Omega$ & $K$ & $\boldsymbol\beta$ & $c$ & $u_\star$ & runs\\
\midrule
Trigonometric (smooth high diffusion) & \texttt{trigonometric} & $\Omega_\square$ & $1$ & $\boldsymbol\beta_0$ & $1$ & $u_s$ & ladder\\
Smooth transport & \texttt{advection\_dominated} & $\Omega_\square$; $\Omega_5$ & $10^{-3}$ & $\boldsymbol\beta_0$ & $1$ & $u_s$ & ladder; stock\\
Variable velocity & \texttt{variable\_velocity} & $\Omega_\square$ & $1$ & $\boldsymbol\beta_v$ & $\tfrac12$ & $u_q$ & ladder\\
Smooth anisotropy & \texttt{anisotropic} & $\Omega_\square$ & $\operatorname{diag}(1,10^{-2})$ & $\boldsymbol\beta_0$ & $1$ & $u_s$ & ladder\\
Oscillatory low diffusion & \texttt{cellular7\_weak} & $\Omega_5$; $\Omega_\square$ & $10^{-5}$ & $\boldsymbol\beta_7$ & $10^{-2}$ & $u_o$ & ladder; stock\\
Closed-loop stress (orthogonal) & \texttt{stress\_main\_orthogonal} & $\Omega_9$ & $K_\perp$ & $\boldsymbol\beta_\perp$ & $10^{-3}$ & $u_9$ & ladder\\
\midrule
Oscillatory source & \texttt{oscillatory\_rhs} & $\Omega_\square$, $\Omega_5$ & $10^{-3}$ & $\boldsymbol\beta_0$ & $1$ & $u_o$ & stock\\
Cellular, four waves & \texttt{cellular4} & $\Omega_\square$ & $10^{-3}$ & $\boldsymbol\beta_4$ & $1$ & $u_o$ & stock\\
Cellular, seven waves & \texttt{cellular7} & $\Omega_\square$ & $10^{-3}$ & $\boldsymbol\beta_7$ & $1$ & $u_o$ & stock\\
Cellular, low diffusion & \texttt{cellular7\_low} & $\Omega_\square$ & $10^{-5}$ & $\boldsymbol\beta_7$ & $1$ & $u_o$ & stock\\
Oscillatory high diffusion & \texttt{cellular7\_high} & $\Omega_\square$, $\Omega_5$ & $1$ & $\boldsymbol\beta_7$ & $10^{-2}$ & $u_o$ & stock\\
Oscillatory directional & \texttt{cellular7\_directional} & $\Omega_\square$, $\Omega_5$ & $K_{1,10^{-3}}$ & $\boldsymbol\beta_7$ & $10^{-2}$ & $u_o$ & stock\\
Weak anisotropy & \texttt{cellular7\_anisotropic} & $\Omega_\square$, $\Omega_5$ & $K_{10^{-4},10^{-6}}$ & $\boldsymbol\beta_7$ & $10^{-2}$ & $u_o$ & stock\\
\bottomrule
\end{tabular}
\caption{Problem classes of Equation~\eqref{eq:adr-h200-problem}. ``Ladder'': the tuned H200
campaign, L1--L7 (stress L4--L7 and L4b); ``stock'': the stock-kernel H200 run on small meshes
(Section~\ref{sec:adr-h200-synth-oscillatory}). A semicolon separates the ladder geometry from the
stock one. Definitions: \texttt{scripts/advection\_diffusion\_reaction/cases/}.}
\label{tab:adr-h200-cases}
\end{table}

\paragraph{Protocol.}
\begin{itemize}
\item \textbf{Fresh} is setup plus the first solve. For the iterative families the setup covers
  the preconditioner (local factorizations, polynomial, AMG hierarchy); for LU it covers CSR
  conversion, analysis and numerical factorization.
\item \textbf{Reused} is the second solve on the same setup, or a solve with the existing LU
  factors.
\item For each system, LU keeps the fastest thread count per metric, chosen among 8--192
  threads; the subscript of P-LU gives it.
\item An iterative solve passes when the solver reports convergence to its internal target
  ($10^{-11}$) and the recomputed true relative residual is at most $10^{-10}$; on variable
  velocity, whose exact solution is quadratic, its $L^2$ error must also stay below $2\cdot10^{-8}$.
  Outcomes are recorded, non-convergence included, and the fastest passing member of each family
  is shown.
\item $^*$ \textbf{FP64 floor.} On the trigonometric and anisotropic classes the round-off floor
  of the FP64 residual, $\varepsilon\||A||x|\|/\|b\|$, grows in proportion to the DOFs: it
  exceeds the internal target from L4--L5 and the acceptance from L7. NC$^*$ (FP64$^*$ in the LU
  tables) marks a run that stopped there with the best solution's accuracy: not a solver failure.
  The graphs show such a run as a hollow marker on a dotted line, at the time its first solve took
  to reach that floor; the grids and tables keep NC$^*$ (Section~\ref{sec:adr-h200-tuned},
  paragraph on the FP64 residual floor).
\end{itemize}

\paragraph{Coverage.} The tuned campaign exported all 114 of its systems (L1--L7 and the stress
mesh L4b) in the layout PyPardiso reads, and all 114 have an LU measurement:
\begin{itemize}
\item 46 from the PARDISO v2 campaign, whose plan came from the unified campaign's exports;
\item 68 from a follow-up run of the same worker on the tuned exports (jobs 26536, 26551 and
  26578), up to L7 at 20 million trace DOFs.
\end{itemize}
The Lmax tier (Section~\ref{sec:adr-h200-synth-lmax}) adds 42 systems of 14--35 million
DOFs without LU. Fresh LU takes 36--53~s at L7, with at most 141~GiB of memory. PyPardiso's 32-bit MKL interface
limits the matrix's own indices (fewer than $2^{31}$ entries), not the factor. Three L7 LU solves
end at a true relative residual of $1.10$--$1.36\cdot10^{-10}$, just above the target: the FP64
floor that the iterative solvers also meet there. Their timings are shown like the others.

\paragraph{Problem sizes.} Table~\ref{tab:adr-h200-levels} gives every level's mesh. L1--L5 and
L4b share their mesh across degrees; L6 and L7 are defined by trace DOFs (8 and 20 million), so
their meshes depend on $p$. Edges are the trace-unknown edges (trace DOFs over $p+1$). The
square mesh serves the trigonometric, transport, variable-velocity and anisotropic cases; the
two annular classes have their own meshes, the stress class from L4.

\begin{table}[ht]
\centering\small
\resizebox{\linewidth}{!}{\input{\ADRResultsPath/h200_tuned/levels_table.tex}}
\caption{Meshes and trace DOFs per level and degree.}
\label{tab:adr-h200-levels}
\end{table}

\subsection{Overview}

Figures~\ref{fig:adr-h200-synth-p6}--\ref{fig:adr-h200-synth-degree-reused} show every family on every
system. Each cell gives the time in ms and the iteration count, with the configuration below.
Shading is the slowdown against the fastest passing family on that system, and a green border
marks the winner.

``NC'' (red) means not passed. ``NC$^*$'' (amber) means stopped at the FP64 residual floor
with the best solution's accuracy (Protocol; 83 runs, all trigonometric or anisotropic, L4--L7 and Lmax;
Table~\ref{tab:adr-h200-outcomes}). A dash marks a family that
was not run on that class. These families were pruned in advance from the Blackwell and unified
results, where they rarely or never converged, or were far slower than the block
preconditioners:
\begin{itemize}
\item BJ+PP on the anisotropic and cellular classes;
\item $p$MG--AMG on the cellular annulus;
\item AMGX and $p$MG--AMG on the closed-loop stress class.
\end{itemize}
A dash in the LU column means no LU measurement (see Coverage).

\begin{figure}[p]
\centering
\includegraphics[width=\linewidth,height=.93\textheight,keepaspectratio]{\ADRResultsFigurePath/h200_synthesis/h200_mesh_p6_fresh.pdf}
\caption{H200 and genoa LU, mesh refinement at $p=6$: fresh solve (setup and first solve). Mesh sizes in Table~\ref{tab:adr-h200-levels}; NC$^*$: FP64 residual floor (Protocol).}
\label{fig:adr-h200-synth-p6}
\end{figure}

\begin{figure}[p]
\centering
\includegraphics[width=\linewidth,height=.93\textheight,keepaspectratio]{\ADRResultsFigurePath/h200_synthesis/h200_mesh_p6_reused.pdf}
\caption{As Figure~\ref{fig:adr-h200-synth-p6}: reused solve (second solve on the same setup; LU with its existing factors).}
\label{fig:adr-h200-synth-p6-reused}
\end{figure}

\begin{figure}[p]
\centering
\includegraphics[width=\linewidth,height=.93\textheight,keepaspectratio]{\ADRResultsFigurePath/h200_synthesis/h200_mesh_p2_fresh.pdf}
\caption{As Figure~\ref{fig:adr-h200-synth-p6}, at $p=2$: fresh solve (setup and first solve). Mesh sizes in Table~\ref{tab:adr-h200-levels}; NC$^*$: FP64 residual floor (Protocol).}
\label{fig:adr-h200-synth-p2}
\end{figure}

\begin{figure}[p]
\centering
\includegraphics[width=\linewidth,height=.93\textheight,keepaspectratio]{\ADRResultsFigurePath/h200_synthesis/h200_mesh_p2_reused.pdf}
\caption{As Figure~\ref{fig:adr-h200-synth-p2}: reused solve (second solve on the same setup; LU with its existing factors).}
\label{fig:adr-h200-synth-p2-reused}
\end{figure}

\begin{figure}[p]
\centering
\includegraphics[width=\linewidth,height=.93\textheight,keepaspectratio]{\ADRResultsFigurePath/h200_synthesis/h200_mesh_p9_fresh.pdf}
\caption{As Figure~\ref{fig:adr-h200-synth-p6}, at $p=9$: fresh solve (setup and first solve). Mesh sizes in Table~\ref{tab:adr-h200-levels}; NC$^*$: FP64 residual floor (Protocol). L7 at $p=9$ exceeds one H200's memory.}
\label{fig:adr-h200-synth-p9}
\end{figure}

\begin{figure}[p]
\centering
\includegraphics[width=\linewidth,height=.93\textheight,keepaspectratio]{\ADRResultsFigurePath/h200_synthesis/h200_mesh_p9_reused.pdf}
\caption{As Figure~\ref{fig:adr-h200-synth-p9}: reused solve (second solve on the same setup; LU with its existing factors).}
\label{fig:adr-h200-synth-p9-reused}
\end{figure}

\begin{figure}[p]
\centering
\includegraphics[width=\linewidth,height=.93\textheight,keepaspectratio]{\ADRResultsFigurePath/h200_synthesis/h200_degree_L4_fresh.pdf}
\caption{H200 and genoa LU, degree sweep at L4 (the same mesh at every degree): fresh solve.
Mesh sizes in Table~\ref{tab:adr-h200-levels}; NC$^*$: FP64 residual floor (Protocol).}
\label{fig:adr-h200-synth-degree}
\end{figure}

\begin{figure}[p]
\centering
\includegraphics[width=\linewidth,height=.93\textheight,keepaspectratio]{\ADRResultsFigurePath/h200_synthesis/h200_degree_L4_reused.pdf}
\caption{As Figure~\ref{fig:adr-h200-synth-degree}: reused solve.}
\label{fig:adr-h200-synth-degree-reused}
\end{figure}


\subsection{Crossovers under refinement and increasing degree}

Figures~\ref{fig:adr-h200-synth-h} and \ref{fig:adr-h200-synth-pfe} follow the smooth cases,
where the families compete most closely. Table~\ref{tab:adr-h200-node-crossover} gives, per
case and degree, the smallest size from which the best iterative solver stays faster than LU.

\begin{figure}[ht]
\centering
\includegraphics[width=\linewidth]{\ADRResultsFigurePath/h200_synthesis/h200_h_crossover.pdf}
\caption{Fresh (top) and reused (bottom) time under refinement at $p=6$; LU points carry their
thread count. Crosses above a panel mark a family that did not pass at that level, one row per
family in its colour. On variable velocity ASM+PP and BJ+PP converge to true residuals below
$10^{-11}$ but miss the case's $L^2$ bound of $2\cdot10^{-8}$ (by 5--53 times at L5--L7; Protocol).
Hollow markers on a dotted line (NC$^*$) give the time a run that did not pass took to reach the
FP64 residual floor with the best accuracy (Protocol).}
\label{fig:adr-h200-synth-h}
\end{figure}

\begin{figure}[ht]
\centering
\includegraphics[width=\linewidth]{\ADRResultsFigurePath/h200_synthesis/h200_p_crossover.pdf}
\caption{As Figure~\ref{fig:adr-h200-synth-h}, against the degree at L4.}
\label{fig:adr-h200-synth-pfe}
\end{figure}

\begin{table}[ht]
\centering\small
\resizebox{\linewidth}{!}{\input{\ADRResultsPath/h200_tuned/node_lu_crossover_table.tex}}
\caption{Iterative versus genoa LU per case and degree, over all 114 H200 systems (L1--L7).
Systems without a passed iterative solve do not count as crossings.}
\label{tab:adr-h200-node-crossover}
\end{table}

\subsection{Hard cases: anisotropy, oscillatory annulus and closed-loop stress}

Figures~\ref{fig:adr-h200-synth-hard6} and \ref{fig:adr-h200-synth-hard9} follow the three cases
where reused LU factors win almost everywhere, now up to L7. The exceptions are the anisotropic
square at L5, $p=6$ and 9, where AMGX's repeated solve takes 0.54 and 0.96 of the reused LU time.
On the cellular annulus LU also keeps the fresh solve from L5 on (1.5--2.2 times faster), and on
stress at L6, $p=2$, and L7. On the anisotropic square at L6--L7 no iterative family passes,
while LU solves every system (at L7, $p=6$, at the residual floor). On the anisotropic square,
ASM+PP(48) and $p$MG--AMG pass up to L4 (L3 at $p=9$). AMGX passes up to L5 and is the only
family left there. On the cellular annulus ASM+PP(24) leads from L2 to L5 and AMGX takes over at
L6--L7, where its iteration count (500--600) is below its counts on the coarse meshes. Table~\ref{tab:adr-h200-largest} lists the
largest passed system of every case and degree, the counterpart of the million-DOF table of the
Blackwell synthesis.

\begin{figure}[ht]
\centering
\includegraphics[width=\linewidth]{\ADRResultsFigurePath/h200_synthesis/h200_hard_cases_p6.pdf}
\caption{Hard cases under refinement at $p=6$: fresh (top) and reused (bottom) time, with the
genoa LU where it was measured. Only ASM+PP and BJ+PP were run on the stress case; its L4b point is
the Blackwell study's second stress mesh (148{,}848 triangles). Crosses and hollow markers as in
Figure~\ref{fig:adr-h200-synth-h}.}
\label{fig:adr-h200-synth-hard6}
\end{figure}

\begin{figure}[ht]
\centering
\includegraphics[width=\linewidth]{\ADRResultsFigurePath/h200_synthesis/h200_hard_cases_p9.pdf}
\caption{As Figure~\ref{fig:adr-h200-synth-hard6}, at $p=9$. L7 does not fit one H200 at this degree; the last point is Lmax (14 million trace DOFs).}
\label{fig:adr-h200-synth-hard9}
\end{figure}

\begin{table}[ht]
\centering
\resizebox{\linewidth}{!}{\input{\ADRResultsPath/h200_tuned/synthesis_largest_table.tex}}
\caption{Largest system per case and degree with a passing iterative solve on the H200: fresh
and reused time in ms (iterations). ``NC'': run, not passed; NC$^*$: stopped at the FP64
residual floor (Protocol); ``--'': not run on that case.}
\label{tab:adr-h200-largest}
\end{table}

\subsubsection{Closed-loop stress}
\label{sec:adr-h200-synth-stress}

The Blackwell study's stress campaign tested three arrangements on a tightly pinched annulus:
trapping, crossing and an orthogonal control (Section~\ref{sec:adr-closed-loop-stress}, where
included). Only the orthogonal control passed there: eight of 60 comparisons, all from ASM+PP,
BJ+PP and $p$MG--AMG, while every AMGX candidate failed. The H200 campaign therefore ran the
orthogonal control only, with the two block preconditioners at PP degree 96:
\begin{itemize}
\item \textbf{Excluded, trapping and crossing:} no candidate converged on Blackwell.
\item \textbf{Excluded, AMGX:} its smoothers converged on no stress system.
\item \textbf{Excluded, $p$MG--AMG:} it converges, but 92--110 times slower than ASM+PP.
\end{itemize}
It extends the control from the Blackwell meshes (L4, 98{,}699 triangles; L4b, 148{,}848) to
L7 at 20 million trace DOFs (Table~\ref{tab:adr-h200-stress}).

\begin{itemize}
\item \textbf{ASM+PP is the stress solver.} It passes every level at $p=2$ and $p=6$ up to L7
  and at $p=9$ up to L5. At $p=6$ its iterations grow from 33 (L4) to 293 (L7), with a reused
  solve of 0.55~s at $10^6$ DOFs and 90~s at $2\cdot10^7$. BJ+PP passes the same systems with
  1.5--9 times as many iterations and is 1.1--6 times slower on reused solves. The one exception
  is $p=2$ at L7, where BJ+PP is 6\,\% faster than ASM+PP. At L6, $p=9$, neither passed.
\item \textbf{Against the Blackwell on the same meshes} ($p=6$): ASM+PP's reused solve takes
  0.55~s on the H200 against 5.20~s on Blackwell at 98{,}699 triangles (9.4 times faster), and
  1.10~s against 12.08~s at 148{,}848 (11.0 times). Its fresh solve is 9.8 and 10.5 times
  faster. The iterations fall from 58 to 33 and from 82 to 45 (true-target trigger). BJ+PP is
  7.8--8.0 times faster on reused solves.
\item \textbf{Against LU, up to L7.} At L4 the genoa LU needs 2.17~s fresh and 0.23~s reused at
  $p=6$, and 0.83~s and 0.05~s at $p=2$; the Blackwell host's LU needed 3.99~s and 0.66~s on the
  $p=6$ system.
  \begin{itemize}
  \item Fresh: the H200 wins at L4, L4b and L5 (the best iterative time is 0.44--0.89 of the
    LU's at L5) and at L6, $p=6$ (0.90). LU wins at L6, $p=2$ (1.2 times faster) and at L7 (1.7
    at $p=6$, 3.1 at $p=2$; 52 and 48~s).
  \item Reused: LU's factors stay faster at every level, 3.8--5.2 times at L5 and 8.5 and 26
    times at L7 (10.6 and 5.6~s against ASM+PP's 90~s at $p=6$ and BJ+PP's 147~s at $p=2$).
  \end{itemize}
\end{itemize}

\begin{table}[ht]
\centering\small
\input{\ADRResultsPath/h200_tuned/stress_table.tex}
\caption{Closed-loop stress, orthogonal control, on the H200 (ms). P-LU rows: the genoa LU on
the same system, fresh/reused thread counts in the subscript.}
\label{tab:adr-h200-stress}
\end{table}

\subsection{Largest mesh per degree}
\label{sec:adr-h200-synth-lmax}

A further tier (level Lmax; \texttt{campaigns/adr\_h200\_tuned/README.md}, finished 30 October
2026) solves every class at $p=2$--9 on the largest mesh one H200 takes at that degree: 34
million trace DOFs at $p=2$ down to 14 million at $p=9$ (Table~\ref{tab:adr-h200-lmax}). $p=6$
was not rerun, because L7 sits at its limit; the degree sweeps take it from L7. This fills the
degrees the L4 sweep left out (Figures~\ref{fig:adr-h200-synth-lmax-grid}--\ref{fig:adr-h200-synth-lmax-hard}).
There is no LU at these sizes: the 42 Lmax systems are listed for it in
\texttt{lu\_followup\_systems.csv}.
\begin{itemize}
\item \textbf{Variable velocity:} $p$MG--AMG passes at every degree in 0.87--2.0~s fresh and
  0.57--0.82~s reused, with 37--40 iterations, 2.6--5 times faster than AMGX (3.3--5.3~s). ASM+PP
  and BJ+PP converge but miss the case's $L^2$ bound (Protocol).
\item \textbf{Transport:} AMGX needs 25--35 iterations, 4.1--8.2~s fresh and 2.5--4.8~s reused;
  $p$MG--AMG takes 19--27~s, ASM+PP 29--81~s and BJ+PP 35--53~s (not passed at $p=2$, 3).
\item \textbf{Trigonometric and anisotropic: no family passes, almost all at the FP64 floor.}
  ASM+PP, AMGX and $p$MG--AMG end at true relative residuals of $1.0$--$8.0\cdot10^{-10}$, the
  floor of 14--34 million DOFs (NC$^*$). AMGX and $p$MG--AMG reach it in about 1.5--14~s fresh,
  ASM+PP in 65--150~s. ASM+PP at $p=8$--9 (trigonometric) and $p=9$ (anisotropic) stagnates at
  the same level but misses the best accuracy; BJ+PP (trigonometric only) stalls for real.
\item \textbf{Cellular annulus:} AMGX passes at $p=3$--9 (68--100~s fresh, 493--862 iterations)
  and diverges at $p=2$; ASM+PP passes only at $p=9$ (139~s, 1{,}805 iterations).
\item \textbf{Closed-loop stress:} ASM+PP passes at every degree, from 417~s (1{,}204 iterations)
  at $p=2$ to 48~s (168) at $p=9$. BJ+PP passes at $p=2$--4, where it is the faster of the two
  at $p=2$ (356~s), and fails from $p=5$.
\item \textbf{The degree on a full GPU.} With the mesh as large as the memory allows, the time
  per solve is nearly flat in $p$ for AMGX on transport and falls with $p$ for $p$MG--AMG on
  variable velocity (2.0 to 0.87~s) and for ASM+PP on stress (from 12 to 3.5~$\mu$s per trace
  DOF).
\end{itemize}

\begin{table}[p]
\centering
\adjustbox{max width=\linewidth,max totalheight=.9\textheight}{\input{\ADRResultsPath/h200_tuned/lmax_table.tex}}
\caption{The largest mesh per degree on one H200 (Lmax). Time in ms (iterations); ``NC'': run,
not passed; NC$^*$: stopped at the FP64 residual floor (Protocol); ``--'': not run on that
class.}
\label{tab:adr-h200-lmax}
\end{table}

\begin{figure}[p]
\centering
\includegraphics[width=\linewidth,height=.93\textheight,keepaspectratio]{\ADRResultsFigurePath/h200_synthesis/h200_degree_lmax_fresh.pdf}
\caption{H200, degree sweep on the largest mesh per degree (Lmax; L7 at $p=6$): fresh solve.
Sizes in Table~\ref{tab:adr-h200-lmax}; LU exists only at $p=6$. NC$^*$: FP64 residual floor
(Protocol).}
\label{fig:adr-h200-synth-lmax-grid}
\end{figure}

\begin{figure}[p]
\centering
\includegraphics[width=\linewidth,height=.93\textheight,keepaspectratio]{\ADRResultsFigurePath/h200_synthesis/h200_degree_lmax_reused.pdf}
\caption{As Figure~\ref{fig:adr-h200-synth-lmax-grid}: reused solve.}
\label{fig:adr-h200-synth-lmax-grid-reused}
\end{figure}

\begin{figure}[ht]
\centering
\includegraphics[width=\linewidth]{\ADRResultsFigurePath/h200_synthesis/h200_p_lmax.pdf}
\caption{Fresh (top) and reused (bottom) time against the degree on the largest mesh per degree,
smooth classes: 7.6 million triangles (34 million trace DOFs) at $p=2$ down to 0.94 million (14
million) at $p=9$; L7 at $p=6$ (1.9 million triangles). Per-degree sizes in
Table~\ref{tab:adr-h200-lmax}. Crosses and hollow markers as in Figure~\ref{fig:adr-h200-synth-h}.}
\label{fig:adr-h200-synth-lmax-smooth}
\end{figure}

\begin{figure}[ht]
\centering
\includegraphics[width=\linewidth]{\ADRResultsFigurePath/h200_synthesis/h200_p_lmax_hard.pdf}
\caption{As Figure~\ref{fig:adr-h200-synth-lmax-smooth}, hard classes: 7.6 to 0.93 million
triangles on the square and the two annuli (Table~\ref{tab:adr-h200-lmax}).}
\label{fig:adr-h200-synth-lmax-hard}
\end{figure}

\subsection{Oscillatory classes with the stock kernels}
\label{sec:adr-h200-synth-oscillatory}

The tuned ladder left out the oscillatory classes of the Blackwell synthesis: high diffusion,
directional and anisotropic cellular coefficients, the oscillatory right-hand side and the other
cellular variants. The stock-kernel H200 run (unified campaign 21958, with the study's own
configurations) solved 37 of the systems that the LU campaign measured on those classes and on
the two classes the ladder has only on the other geometry. All of them are small: original
meshes, L1 and L2, at most 85{,}120 trace DOFs, $p=1$--6. Table~\ref{tab:adr-h200-oscillatory}
gives the largest system per class.
\begin{itemize}
\item \textbf{Reused LU factors win.} The best repeated stock solve beats them on 2 of the 37
  systems only, the annular transport and oscillatory right-hand-side systems (9.6 against 11.1
  and 11.5~ms). Elsewhere the repeated iterative solve takes a median 6.9 times the reused LU
  time, up to 41 times.
\item \textbf{Fresh solves split by class.} The GPU wins 12 of 37:
  \begin{itemize}
  \item 7 of the 12 high-diffusion systems ($p$MG--AMG);
  \item the annular transport, oscillatory right-hand-side and anisotropic cellular systems;
  \item one of the six directional annulus systems;
  \item one of the two anisotropic cellular squares.
  \end{itemize}
  LU keeps the fresh solve on the cellular squares (low diffusion, the original cellular cases)
  and on the directional square, where ASM+PP needs 2.0 times the LU time at 85{,}120 DOFs.
\item \textbf{Stock kernels at small sizes, so no crossover.} Below $10^5$ DOFs the tuned kernels
  gain only 1.0--3.4 times (Section~\ref{sec:adr-h200-tuned}), so these results would not change
  much with them. Where these classes cross LU needs the tuned kernels at L3--L5; the ladder did
  not run them.
\end{itemize}

\begin{table}[ht]
\centering\scriptsize
\resizebox{\linewidth}{!}{\input{\ADRResultsPath/h200_tuned/stock_oscillatory_table.tex}}
\caption{Classes outside the tuned ladder, stock kernels on the H200 against the genoa LU:
largest system per class with both (ms). Fresh: setup and first solve; repeated: second solve on
the same setup. Families: fastest fresh / fastest repeated.}
\label{tab:adr-h200-oscillatory}
\end{table}

\subsection{End to end: assembly and one solve}
\label{sec:adr-h200-synth-e2e}

The comparisons above start from an assembled matrix. Table~\ref{tab:adr-h200-e2e} adds the
assembly, on the 24 systems that both the H200 and hybridge's host pipeline on a genoa node ran
(the \texttt{cpu\_adr\_scaling} campaign, same plan keys).
\begin{itemize}
\item \textbf{genoa:} the Numba assembly of the host pipeline (its warm run, at the fastest of
  1--192 threads) plus the LU campaign's fresh solve of the same system (analysis, factorization
  and first solve). The pipeline's own warm solve is not used: PyPardiso skips the factorization
  when it is handed the matrix it factored last, so that solve reused the cold run's factors.
  Reconstruction and postprocessing come on top on both machines and are left out.
\item \textbf{H200:} the device assembly of the system job (one cold assembly in the solver
  process), plus the fastest passed fresh iterative solve.
\end{itemize}
Findings:
\begin{itemize}
\item \textbf{From L3 the H200 is faster end to end on every shared system}, cellular annulus
  included.
  \begin{itemize}
  \item L4: 5.8 times on trigonometric, 5.7 on variable velocity, 3.3 on anisotropic and 3.0 on
    stress (where the host assembly alone takes 0.95~s).
  \item L3: 4.1 on transport, 3.5 on variable velocity, 2.0 on trigonometric and anisotropic,
    1.3 on the cellular annulus.
  \end{itemize}
\item \textbf{At L2, $p=6$, the two are about even} (1.0--1.4 times), except on the cellular
  annulus, where genoa needs 0.58 of the H200 time.
\item \textbf{At L1, and at L2 with $p=2$, genoa is faster} (0.20--0.55 of the H200 time). Here
  the H200's cold assembly (143--374~ms) costs more than genoa's assembly and LU together.
\end{itemize}

\begin{table}[ht]
\centering\scriptsize
\input{\ADRResultsPath/h200_tuned/end_to_end_table.tex}
\caption{Assembly plus one fresh solve (ms). genoa: host-pipeline assembly at its fastest thread
count plus the LU campaign's fresh solve (subscripts: threads); reconstruction excluded. H200:
device assembly plus the fastest passed fresh iterative solve. Last column: genoa time over H200
time; above one the H200 is faster.}
\label{tab:adr-h200-e2e}
\end{table}

\subsection{Practical interpretation}

\begin{itemize}
\item \textbf{One solve per matrix: use the GPU, except on the hard cases at scale.} Setup plus
  the first solve beats LU's analysis and factorization on 89 of the 104 systems with a passed
  iterative solve. On the smooth cases the GPU needs 2--18\,\% of the LU time at L5--L7 (36--53~s
  of LU at L7). LU keeps the fresh solve on the smallest meshes (L1 at $p=2$, cellular at L1--L2),
  on the cellular annulus from L5 (1.5--2.2 times faster) and on closed-loop stress at L6, $p=2$,
  and L7 (1.2--3.1 times).
\item \textbf{Many solves with one matrix: smooth cases go to the GPU, hard cases to LU.} A
  reused LU solve takes about 3~ms at $3.6\cdot10^4$ DOFs, 0.23--0.25~s at $10^6$ and 5--11~s at
  $2\cdot10^7$. The GPU wins 42 of the 104 repeated solves:
  \begin{itemize}
  \item transport (15 of 20), variable velocity (14 of 20) and trigonometric (11 of 18), from
    $3\cdot10^4$--$3.4\cdot10^5$ DOFs at $p=6$ and 9;
  \item at $p=2$ these cases cross only from about $4\cdot10^5$ (variable velocity) or 8 million
    DOFs (transport, trigonometric);
  \item the anisotropic square only at L5, $p=6$ and 9.
  \end{itemize}
  At the largest size where an iterative solve passed, the repeated solve takes 0.06--0.10 of
  the reused LU time on variable velocity ($p$MG--AMG, up to 20 million DOFs), 0.29--0.36 on
  trigonometric (AMGX, 8 million) and 0.33--0.71 on transport (AMGX, 8--20 million). On the cellular annulus and on
  closed-loop stress the reused factors stay faster at every size up to L7 (3.5--35 and 1.3--26
  times). On the anisotropic square at L6--L7 LU is the only solver that meets the target.
\item \textbf{Which GPU family at scale.}
  \begin{itemize}
  \item AMGX keeps 15--27 iterations up to L7. It wins the trigonometric case at L5--L6, the
    anisotropic case at L5, transport at L6--L7 and the cellular annulus at L6--L7, at every
    degree it was run.
  \item $p$MG--AMG wins variable velocity at every level from L5 (and from L3 at $p=6$).
  \item ASM+PP wins transport and the cellular annulus at L5 and the closed-loop stress case at
    L5--L7, where only the block preconditioners were run; BJ+PP takes stress at L7, $p=2$.
  \end{itemize}
  LU's memory is not the limit on a genoa node (at most 141~GiB at L7, out of 754~GiB). Its
  time is: 36--53~s fresh at L7, against 1.0--1.2~s for $p$MG--AMG on variable velocity and
  5.2--5.5~s for AMGX on transport.
\item \textbf{The largest problem on one H200} (34 million trace DOFs at $p=2$, 14 million at
  $p=9$): the smooth classes solve in seconds ($p$MG--AMG 0.9--2.0~s on variable velocity, AMGX
  4--8~s on transport), the hard ones in minutes (AMGX 68--100~s on the cellular annulus, ASM+PP
  48--417~s on stress). On the trigonometric and anisotropic classes no solver can meet the fixed
  residual target there; the FP64 floor is above it.
\item \textbf{Compared with the Blackwell workstation.} The genoa LU is 1.1--1.9 times faster
  fresh and 1.5--3.2 times faster reused than the Blackwell host's LU, and the H200 iterative
  solves are up to 8 times faster than the Blackwell's (Section~\ref{sec:adr-h200-tuned}). The
  net effect: the fresh-solve picture moves further towards the GPU, while the reused-solve
  crossovers sit at somewhat larger meshes than against the Blackwell host LU.
\end{itemize}
